\subsection{Applications: II. Vector spaces of continuous complex functions}

Let X be a nonempty compact space, H a linear subspace of the complex Banach space \(\mathcal{C}(\mathrm{X};\mathbf{C})\) that contains the constants and separates the points of X. The set of real parts \(\mathcal{R}(f)\) of the functions \(f \in H\) is a linear subspace \(H_{r}\) of the real vector space \(\mathcal{C}(\mathrm{X};\mathbf{R})\) ; for every \(f \in H\) , the set \(H_{r}\) also contains \(\mathcal{I}(f) = \mathcal{R}(-if)\) ; it follows that \(H_{r}\) separates the points of X, because the relation \(h(x) = h(y)\) for all \(h \in H_{r}\) implies that \(\mathcal{R}(f(x)) = \mathcal{R}(f(y))\) and \(\mathcal{I}(f(x)) = \mathcal{I}(f(y))\) and so \(f(x) = f(y)\) for all \(f \in H\) . The \(H_{r}\) -extremal points in X are again called H-extremal, the set of them is denoted \(\operatorname{Ch}_{\mathcal{H}}(\mathrm{X})\) , and the closure of the latter set is denoted \(\check{\operatorname{S}}_{\mathcal{H}}(\mathrm{X})\) . The analogues of Props. 5 and 7 are the following:

PROPOSITION 9. --- For every function \(f \in \mathcal{H}\) , \(\mathrm{Ch}_{\mathcal{H}}(\mathrm{X})\) intersects the set of points where \(|f|\) attains its supremum.

We may limit ourselves to the case that \(f\) is not the constant 0. Let \(a\) be a point of \(X\) where \(|f|\) attains its supremum, and set \(g = f / f(a)\) ; then \(g(a) = 1\) and \(|g(x)| \leqslant 1\) for all \(x \in X\) , whence

\[
\mathscr {R} (g (a)) = 1 \quad \text { and } \quad \mathscr {R} (g (x)) \leqslant 1 \quad \text { for   all } x \in X.
\]

By Prop. 5 of No.~3 applied to \(\mathcal{H}_r\) , there exists \(b \in \mathrm{Ch}_{\mathcal{H}}(\mathrm{X})\) where \(\mathcal{R}(g(x))\) attains its supremum 1, whence \(|g(b)| = 1\) since \(|g(b)| \leqslant 1\) ; it follows that \(|f(b)| = |f(a)| \geqslant |f(x)|\) for all \(x \in \mathrm{X}\) .

PROPOSITION 10. --- Let F be a closed subset of X. The following properties are equivalent:

\begin{enumerate}
\def\labelenumi{\alph{enumi})}
\item
  F contains \(\check{\mathrm{S}}_{\mathcal{H}}(\mathrm{X})\)
\item
  For every function \(f \in \mathcal{H}\) , \(F\) intersects the set of points of \(X\) where \(|f|\) attains its supremum.
\item
  For every point \(x \in X\) , there exists a positive measure \(\mu\) of total mass 1 on X such that \(\operatorname{Supp}(\mu) \subset \operatorname{F}\) and \(f(x) = \int f \, d\mu\) for every function \(f \in H\) .
\end{enumerate}

Let us prove the equivalence of the conditions \(a)\) and \(c)\) : let \(f = f_{1} + if_{2}\) with \(f_{1}, f_{2}\) in \(\mathcal{H}_{r}\) ; the relation \(f(x) = \int f d\mu\) is equivalent to the two relations \(f_{1}(x) = \int f_{1} d\mu\) and \(f_{2}(x) = \int f_{2} d\mu\) ; it thus suffices to apply to \(\mathcal{H}_{r}\) the equivalence of the conditions \(a)\) and \(c)\) of Prop. 7 of No.~3. The fact that \(a)\) implies \(b)\) follows from Prop. 9. Let us show that \(b)\) implies \(a)\) ; this is a matter of seeing that if \(b)\) is verified, then, for every \(h \in \mathcal{H}_{r}\) , \(F\) intersects the set of points where \(h\) attains its infimum in \(X\) . The condition \(b)\) implies that \(F\) is nonempty; since \(F\) is compact, there exists \(a \in F\) such that \(h(a) \leqslant h(y)\) for all \(y \in F\) . Let \(f \in \mathcal{H}\) be such that \(h = \mathcal{R}(f)\) ; for every \(\varepsilon > 0\) , the function \(g = f - h(a) + \varepsilon\) belongs to \(\mathcal{H}\) , and

\[
\mathcal {R} (g (y)) = h (y) - h (a) + \varepsilon \geqslant \varepsilon
\]

for all \(y \in F\) . Let c be the supremum of \(|g|\) in X, and set \(b = c^{2}/2\varepsilon\) ; for every \(y \in F\) ,

\[
\left| g (y) - b \right| ^ {2} = \left| g (y) \right| ^ {2} - 2 b \mathcal {R} (g (y)) + b ^ {2} \leqslant c ^ {2} - 2 b \varepsilon + b ^ {2} = b ^ {2},
\]

in other words, the supremum in F of the function \(|g - b|\) is \(\leqslant b\) . Since \(g - b \in \mathcal{H}\) , the hypothesis on F implies that \(|g - b| \leqslant b\) , whence

\[
b ^ {2} \geqslant | g - b | ^ {2} = | g | ^ {2} - 2 b \mathcal {R} (g) + b ^ {2}
\]

and so \(\mathcal{R}(g) \geqslant |g|^{2}/2b \geqslant 0\) ; since \(\mathcal{R}(g) = h - h(a) + \varepsilon\) , and \(\varepsilon > 0\) is arbitrary, we have \(h \geqslant h(a)\) , and \(h(a)\) is the infimum of h in X, which completes the proof.

Remark. --- If f is a continuous real function, a point where \(|f|\) attains its supremum is a point where one of the functions f, -f attains its supremum. For a vector space H of continuous real functions satisfying the hypotheses of No.~3, the Props. 9 and 10 are thus trivial corollaries of Props. 5 and 7, respectively.

